\chapter*{Appendices}
\addcontentsline{toc}{chapter}{Appendices}

\section*{Appendix A: Example Learner Output (raw JSON)}
\label{app:json}

The listing below is the verbatim structured output produced by the learner VLM
(\texttt{qwen3.6-plus}) on the \texttt{wrong\_color} clip in the blind v3 run. For each
guiding question the model returns an answer and a short visual justification, then the
overall verdict, anomaly type, score, confidence and a one-sentence reason. This is the
per-question evidence that makes the decision explainable, summarised in
Table~\ref{tab:example}.

\footnotesize
\begin{verbatim}
{
  "question_answers": [
    { "question_id": 1, "answer": "YES",
      "evidence": "At timestamp 00:10, the board is fully visible and clearly
                   holds three pens in a row." },
    { "question_id": 2, "answer": "YES",
      "evidence": "All three pens are oriented vertically with the caps
                   pointing downwards." },
    { "question_id": 3, "answer": "NO",
      "evidence": "The leftmost pen has a blue barrel, while the middle and
                   right pens have black/dark grey barrels." },
    { "question_id": 4, "answer": "YES",
      "evidence": "All pens appear to have caps on the bottom and are of
                   similar length." },
    { "question_id": 5, "answer": "YES",
      "evidence": "The board moves continuously from left to right throughout
                   the clip without stopping." },
    { "question_id": 6, "answer": "YES",
      "evidence": "The pens remain fixed in their relative positions on the
                   board as it slides." }
  ],
  "verdict": "ANOMALY",
  "anomaly_type": "wrong_color",
  "anomaly_score": 1.0,
  "confidence": "HIGH",
  "reason": "The video shows three pens on the board, but one is blue while the
             other two are black, violating the requirement that all pens be
             identical in color."
}
\end{verbatim}
\normalsize
